% Options for packages loaded elsewhere
\PassOptionsToPackage{unicode}{hyperref}
\PassOptionsToPackage{hyphens}{url}
\PassOptionsToPackage{dvipsnames,svgnames,x11names}{xcolor}
\documentclass[
  11pt,
  a4paper,
]{article}
\usepackage{xcolor}
\usepackage[margin=24mm]{geometry}
\usepackage{amsmath,amssymb}
\setcounter{secnumdepth}{-\maxdimen} % remove section numbering
\usepackage{iftex}
\ifPDFTeX
  \usepackage[T1]{fontenc}
  \usepackage[utf8]{inputenc}
  \usepackage{textcomp} % provide euro and other symbols
\else % if luatex or xetex
  \usepackage{unicode-math} % this also loads fontspec
  \defaultfontfeatures{Scale=MatchLowercase}
  \defaultfontfeatures[\rmfamily]{Ligatures=TeX,Scale=1}
\fi
\usepackage{lmodern}
\ifPDFTeX\else
  % xetex/luatex font selection
\fi
% Use upquote if available, for straight quotes in verbatim environments
\IfFileExists{upquote.sty}{\usepackage{upquote}}{}
\IfFileExists{microtype.sty}{% use microtype if available
  \usepackage[]{microtype}
  \UseMicrotypeSet[protrusion]{basicmath} % disable protrusion for tt fonts
}{}
\makeatletter
\@ifundefined{KOMAClassName}{% if non-KOMA class
  \IfFileExists{parskip.sty}{%
    \usepackage{parskip}
  }{% else
    \setlength{\parindent}{0pt}
    \setlength{\parskip}{6pt plus 2pt minus 1pt}}
}{% if KOMA class
  \KOMAoptions{parskip=half}}
\makeatother
\usepackage{longtable,booktabs,array}
\newcounter{none} % for unnumbered tables
\usepackage{calc} % for calculating minipage widths
% Correct order of tables after \paragraph or \subparagraph
\usepackage{etoolbox}
\makeatletter
\patchcmd\longtable{\par}{\if@noskipsec\mbox{}\fi\par}{}{}
\makeatother
% Allow footnotes in longtable head/foot
\IfFileExists{footnotehyper.sty}{\usepackage{footnotehyper}}{\usepackage{footnote}}
\makesavenoteenv{longtable}
\setlength{\emergencystretch}{3em} % prevent overfull lines
\providecommand{\tightlist}{%
  \setlength{\itemsep}{0pt}\setlength{\parskip}{0pt}}
% Shared layout only. Research content is generated from protocol_EN_full.md.
% "Protocol v<version>" must match 01_protocol/project_settings.json (checked by build_protocol.py).
\usepackage{microtype}
\usepackage{fancyhdr}
\usepackage{xurl}
\setlength{\emergencystretch}{3em}
\pagestyle{fancy}
\fancyhf{}
\fancyhead[L]{\small Open-window precision biomarkers}
\fancyhead[R]{\small Protocol v3.0}
\fancyfoot[L]{\small Pre-archive working protocol}
\fancyfoot[R]{\thepage}
\setlength{\headheight}{14pt}
\usepackage{bookmark}
\IfFileExists{xurl.sty}{\usepackage{xurl}}{} % add URL line breaks if available
\urlstyle{same}
\hypersetup{
  pdftitle={Precision biomarkers of the exercise-induced immune “open window”: a scoping review of the validation readiness of conventional and omics candidates},
  colorlinks=true,
  linkcolor={Maroon},
  filecolor={Maroon},
  citecolor={Blue},
  urlcolor={teal},
  pdfcreator={LaTeX via pandoc}}

\author{}
\date{}

\begin{document}

\section{Precision biomarkers of the exercise-induced immune ``open
window'': a scoping review of the validation readiness of conventional
and omics
candidates}\label{precision-biomarkers-of-the-exercise-induced-immune-open-window-a-scoping-review-of-the-validation-readiness-of-conventional-and-omics-candidates}

\textbf{Protocol version:} 3.0 (draft for team freeze) --- 4 October
2026\\
\textbf{Status:} Pre-archive working protocol; formal searches,
screening and synthesis have not begun.\\
\textbf{Target journal:} \emph{Frontiers in Immunology}, under the
Systematic Review article type, which includes scoping reviews;
specialty section to be selected. Submission will include the PRISMA-ScR
checklist and flow diagram as supplementary material and an AI-use
disclosure in both Methods and Acknowledgments stating tool name,
version, model and source.\\
\textbf{Registration and archiving:} Public GitHub repository with a
GitHub Release archived by Zenodo, which assigns a dated, immutable
version DOI and a concept DOI on release. No release exists, no DOI has
been assigned, and this project has never been registered in any
registry; PROSPERO does not accept scoping reviews and OSF Registries is
no longer the planned route.\\
\textbf{Team:} Access to all five databases and a second human reviewer
are confirmed. Names, adjudicator, search peer reviewer, funding and
competing-interest declarations remain to be supplied.

This is the canonical English protocol. The Chinese protocol is an
aligned implementation version. The rendered English source and PDF are
generated from this complete document; they are not a shortened
independent protocol. Operational appendices in
\texttt{03\_search}--\texttt{06\_synthesis} implement the decisions
below. This version follows three pre-archive activities disclosed in J:
an internal protocol (version 1.0, January 2026), public exploratory
reconnaissance and source verification on 2 October 2026 (Codex), and a
collision re-check of the narrowed framing against published and
registered reviews on 4 October 2026 (Claude Code;
\texttt{02\_preliminary/collision\_recheck\_2026-10-04.md}). The
protocol date is not a completed formal search cutoff. No search yield,
eligibility decision, included-study count or review result is claimed
here.

\subsection{A. Background and
rationale}\label{a.-background-and-rationale}

Exercise can alter circulating immune-cell counts, immune mediators,
mucosal measurements and immune-cell function over different time
scales. The ``open window'' has been used to describe the hours to days
after strenuous or prolonged exercise during which immune changes have
been proposed to coincide with reduced host protection {[}15,16{]}.
Whether such a period exists remains debated: lymphocytopenia may
reflect redistribution rather than loss of whole-body function, and
plasma-volume change, circadian timing and sampling alter interpretation
{[}1,2{]}. This review does not adjudicate that debate. It uses the open
window solely as an organizing construct that defines the post-exercise
period (0--72 hours after cessation of an identifiable bout), the
sampling time points and the marker families within which candidates are
charted, irrespective of whether original authors endorse, reject or
omit the hypothesis.

Candidates are charted by marker family. Conventional families include
leukocyte and lymphocyte redistribution, natural-killer-cell counts and
cytotoxicity, neutrophil function, T-cell subsets and function, mucosal
secretory IgA, cytokines and other inflammatory mediators, and
immunoglobulins and complement. Omics families include immune-cell
transcriptomics, proteomics, metabolomics and lipidomics, epigenomics,
non-coding RNA and single-cell measurements. Higher measurement
dimensionality does not alone confer immune specificity or individual
predictive value. A group-average response, an in-silico pathway
annotation, an association with symptoms and a validated individual
prediction address different questions. Candidate usefulness also
depends on the exercise exposure, tissue or fluid, analytical method,
sampling interval and intended monitoring use.

The intended contribution is a per-marker \textbf{validation-readiness
map} of candidates measured within the organizing window. Throughout
this protocol and its appendices, a validation-readiness map means a
descriptive, domain-by-domain profile, for each marker or marker family,
of which of the seven precision-evidence domains defined in C have
actually been investigated and with what result. It is not a score,
ranking, grade or development ladder, and every later prohibition on
readiness scores or rankings refers to this definition. The map
emphasizes individual baselines, within-person variation, repeatability,
time-resolved interpretation, functional or clinical links and
independent validation. It is not a catalogue that labels every
statistically significant omics feature a precise biomarker. Analytical
and clinical validation must be interpreted for a specified context of
use {[}3{]}.

In their 2017 review of immune recovery after exercise, Peake and
colleagues concluded that, ``except for salivary IgA, clear and
consistent markers of this immunodepression remain elusive'' {[}15{]};
that statement is the auditable starting point for this map. The debate
on whether the open window exists is held by Peake et al.~2017 {[}15{]},
Campbell and Turner 2018 {[}1{]}, Simpson et al.~2020 {[}16{]}, Reitzner
and Brodin 2026 {[}19{]} and Xie et al.~2026 {[}17{]}; this review does
not join it. Narrative reviews have also proposed marker agendas around
the open window: Xie et al.~2026 in \emph{Frontiers in Immunology}
(tissue-resident memory T cells and their circulating precursors)
{[}17{]}; Shi et al.~2025, also in \emph{Frontiers in Immunology}
(workload-dependent recovery of conventional indices) {[}18{]}; Reitzner
and Brodin 2026 (the immune system as a systems-level biosensor of
health and performance) {[}19{]}; and Mănescu et al.~2026 (molecular
biomarkers of training responses) {[}20{]}. Davison et al.~2025
questioned the clinical relevance of ex vivo immune assays {[}2{]}. None
of these charts human evidence marker by marker against prespecified
validation domains with a reproducible search, and systematic reviews to
date address single marker families or single populations:
exercise-associated non-coding RNAs, B-cell-related immunity, TH1/TH2
markers in older adults and exercise immunometabolomics {[}4--7{]}. The
registered reviews nearest to this protocol differ in
Population--Concept--Context: PROSPERO CRD42020186264 is confined to
salivary IgA and respiratory-infection risk; PROSPERO CRD42022324139 is
confined to the acute effect of high-intensity interval training; and
OSF osf.io/4hsdr shares the JBI and PRISMA-ScR method stack but is
restricted to athletes aged 10--17 years and to conventional markers,
without omics, open-window or validation dimensions, whereas this review
includes adults only. The 4 October 2026 re-check that identified these
items is filed at
\texttt{02\_preliminary/collision\_recheck\_2026-10-04.md}. We make no
``first review'' or ``no competing review'' claim. The proposed
distinction is the joint, time-resolved charting of conventional and
omics candidates within one organizing window against separately
recorded precision domains. That distinction will be rechecked against
published and registered reviews before the v3.0 archive release and at
the final search update.

A scoping review is appropriate because the primary questions concern
the distribution, measurement, interpretation and validation of
heterogeneous evidence, rather than one intervention effect or
diagnostic threshold. We will follow JBI scoping-review methods {[}8{]},
report with PRISMA-ScR {[}9{]}, and report searches with PRISMA-S
{[}10{]}. Prespecified, design-appropriate critical appraisal will
support interpretation. We will not conduct a meta-analysis or turn
appraisal into a clinical ranking.

\subsection{B. Objective and review
questions}\label{b.-objective-and-review-questions}

This scoping review will map human evidence on candidate immune
biomarkers measured within the exercise-induced ``open window'', used
here solely as an organizing construct that defines the post-exercise
period after an identifiable bout, the sampling time points and the
marker families to be charted. Across conventional immune indices and
omics candidates in generally healthy adults after single identifiable
bouts (Core A) and across repeated bouts in the same individuals
(Support B), we will chart, marker by marker, which of seven
precision-evidence domains have actually been investigated: analytical
reliability, temporal validity, immune specificity, individualization,
functional or clinical linkage, independent validation and use, and
feasibility. The resulting per-marker validation-readiness map is a
descriptive domain-by-domain profile, not a score, ranking or
development ladder. The review does not test whether an open window of
increased infection susceptibility exists, and original authors'
endorsement or rejection of that hypothesis neither determines
eligibility nor weights evidence.

\begin{enumerate}
\def\labelenumi{\arabic{enumi}.}
\tightlist
\item
  \textbf{Measurements and context:} Which marker families, individual
  biomarkers, analytical platforms, immune compartments and sample
  matrices have been measured within the organizing window (0--72 hours
  after cessation), in which adult populations and exercise settings?
\item
  \textbf{Timing, confounding and interpretation:} Which time points
  within the window (immediate, 30 minutes to \textless3 hours, 3 to
  \textless24 hours, 24 to 72 hours) and which outside-window recovery
  points are covered; how are cell composition, hydration, sampling and
  other relevant confounders addressed; and how do original authors
  interpret their findings in relation to the open window (charted
  descriptively only)?
\item
  \textbf{Individualization:} What evidence exists on personal
  baselines, within-person variability, repeatability across bouts,
  response heterogeneity, subgroup interpretation or individualized
  models, and how has it been validated?
\item
  \textbf{Function, outcomes and use:} Which functions or clinical
  outcomes were merely co-measured, which were associated with
  biomarkers, and which monitoring or prediction uses underwent
  independent validation? Together with questions 1--3 this yields the
  per-marker validation-readiness map and the development requirements
  that remain.
\end{enumerate}

Positive, null and inconsistent findings are all relevant. The questions
do not presuppose that exercise is immunosuppressive, nor test whether
an open window exists, nor assume that any candidate is ready to predict
infection.

\subsection{C. Framework and operational
definitions}\label{c.-framework-and-operational-definitions}

Eligibility uses Population--Concept--Context (PCC). ``Post-exercise''
begins at cessation of an identifiable bout or competition. Exact times
and their reported origin are retained. During-exercise samples can be
ancillary observations in an eligible study but cannot alone establish
post-exercise eligibility.

The \textbf{organizing post-exercise window} is the period from
cessation to 72 hours inclusive. It is a charting frame, not an
eligibility cap: a study qualifies through at least one post-cessation
immune-related sample, and samples taken after 72 hours in an otherwise
eligible study are extracted in full and tagged ``outside-window
recovery'' because they document recovery endpoints. No exclusion code
refers to the window. \textbf{Author open-window interpretation} is a
descriptive tag recording whether the original authors used the
open-window term, invoked immunosuppression or susceptibility without
the term, interpreted changes as redistribution, or offered no window
interpretation; it affects neither eligibility nor weighting.

For this project, \textbf{precision monitoring} means interpreting an
immune-related measurement within the organizing window for a defined
person or subgroup, exercise context, sampling time and intended
monitoring use. Seven domains are recorded separately:

{\def\LTcaptype{none} % do not increment counter
\begin{longtable}[]{@{}
  >{\raggedright\arraybackslash}p{(\linewidth - 2\tabcolsep) * \real{0.5000}}
  >{\raggedright\arraybackslash}p{(\linewidth - 2\tabcolsep) * \real{0.5000}}@{}}
\toprule\noalign{}
\begin{minipage}[b]{\linewidth}\raggedright
Domain
\end{minipage} & \begin{minipage}[b]{\linewidth}\raggedright
Evidence to chart
\end{minipage} \\
\midrule\noalign{}
\endhead
\bottomrule\noalign{}
\endlastfoot
Analytical reliability & Assay performance, limits, quality control,
reproducibility and pre-analytical handling \\
Temporal validity & Exact exercise-to-sample intervals, position
relative to the organizing window, longitudinal coverage and how
response or recovery is defined \\
Immune specificity & Direct immune measurement, cellular origin,
composition adjustment and relevant alternative explanations \\
Individualization & Personal baselines, within-person noise,
repeated-bout reproducibility, response heterogeneity, subgroup or
individual models \\
Functional or clinical linkage & Direct immune-function measurements,
tested associations and explicitly defined clinical outcomes \\
Independent validation and use & Technical confirmation versus
independent participants/cohorts, model validation and tested decision
use \\
Feasibility & Sampling burden, accessibility, cost, turnaround and field
implementation, when reported \\
\end{longtable}
}

These are project operational domains, not a validated scale, regulatory
designation or ordinal development ladder. There will be no combined
``precision score''; domain evidence is summarized only as the
per-marker validation-readiness map defined in A. Omics measurement,
machine learning or a group mean difference alone does not establish
personalized usefulness. Individualization or clinical validation is
\textbf{not required for inclusion}; absence of such reporting is itself
a mapping finding.

A \textbf{response} is an observed post-exercise measurement or
contrast. A \textbf{recovery trajectory} requires longitudinal
observations and an explicit endpoint; a single early sample or a
nonsignificant comparison with baseline does not establish recovery.
Cell-count normalization, biomarker normalization, immune-function
recovery and clinical recovery remain distinct.

\textbf{Immune linkage} requires a direct
immune-cell/count/subset/effector measure; a recognized immune mediator
such as an immunoglobulin, cytokine or complement component; omics
measured in immune cells; or biofluid omics with an immune objective or
analysis stated in the original aims/methods and identifiable immune
candidates. Original-study preregistration is not required. Mere
post-hoc immune enrichment in an otherwise general muscle or metabolic
study is insufficient for the main map. Cortisol, creatine kinase,
lactate or heart-rate variability alone do not qualify, but may be
charted as covariates in otherwise eligible studies.

Direct functional assays are distinguished from computational
annotation. Co-measurement of a function or clinical outcome does not
establish a biomarker--outcome association. Self-reported symptoms,
clinician-defined syndromes and pathogen-confirmed infection are
separate outcome classes. Association, temporal prediction, calibration,
discrimination and clinical utility are recorded separately. We use
``prediction'' for future outcome prediction here, without assuming the
treatment-response meaning of a predictive biomarker in other
taxonomies.

Missing information is coded \textbf{NR} (not reported), \textbf{NA}
(not applicable) or \textbf{UNCLEAR} (information present but
insufficient or ambiguous after checking). NR is not evidence of
biological absence or methodological failure.

\subsection{D. Eligibility criteria}\label{d.-eligibility-criteria}

\subsubsection{D.1 Population}\label{d.1-population}

Eligible participants are adults aged 18 years or older, of any sex or
training status, who are generally healthy at baseline and do not have
an active infection. Athletes, recreational exercisers and untrained
participants are eligible; training strata are charted separately. Older
adults are eligible and flagged as a subgroup. Children and adolescents
(\textless18 years) are excluded; mixed-age cohorts are eligible only
where an adult stratum is separable in the report. This boundary
distinguishes the review from the registered youth-athlete scoping
review osf.io/4hsdr (10--17 years). Extractable healthy adult strata
within cohorts of mixed health status are also eligible. Cohorts with
known clinical/metabolic disease or disease-treatment exercise
programmes are excluded from the main map. Ambiguous metabolic-health
descriptions, including insulin resistance without a clear clinical
classification, require explicit adjudication using the reported
diagnosis and separability of strata, rather than an automatic healthy
designation.

\subsubsection{D.2 Exercise context and evidence
streams}\label{d.2-exercise-context-and-evidence-streams}

\textbf{Core A --- acute challenge:} An identifiable single bout or
competition of endurance, resistance, interval/sprint or mixed exercise,
with at least one post-cessation immune-related sample and either a
pre-bout baseline or a suitable non-exercise/comparator condition. No
intensity, duration or mode threshold is imposed: any identifiable bout
qualifies, and intensity, duration and mode are recorded as reported and
used as evidence-map strata, never as eligibility filters, because the
open-window literature defines ``prolonged'' and ``intense''
inconsistently and a threshold would discard comparator evidence.
Immediate samples are eligible and there is no minimum 30-minute
follow-up. The organizing window of 0--72 hours after cessation frames
the charting; it is not a maximum interval, and samples beyond 72 hours
in an eligible study are extracted and tagged ``outside-window
recovery'' (C). During-only studies are ineligible. Repeated acute
challenges before and after a training programme qualify for A; training
moderation is extracted separately. Immediate-only studies contribute
response evidence, not a complete recovery trajectory.

\textbf{Support B --- repeated post-exercise windows in the same
people:} The same immune-related marker is measured in the same
individuals in the post-exercise windows after at least two identifiable
bouts or competitions, with a known or reliably recoverable
exercise-to-sampling relation and a baseline or appropriate comparator
enabling interpretation. B is the cross-bout monitoring layer and the
main evidence source for the individualization and repeatability
domains; it is presented separately from A. A cohort may meet both A and
B; linked cohort identifiers prevent duplicate overall counting. General
weekly or seasonal resting surveillance with unknown last-bout timing,
training pre/post snapshots with unknown exercise-to-sampling intervals,
and habitual-exercise cross-sectional comparisons are background only. A
long observation period cannot substitute for an identifiable
post-exercise time relation.

Studies involving supplementation, diet or other co-interventions
qualify only where an exercise contrast is interpretable separately, for
example through an appropriate control/placebo arm or factorial
contrast. Otherwise the inseparable intervention is ineligible. Prior
training, co-interventions and recovery arrangements are extracted.

\subsubsection{D.3 Concept, design and publication
status}\label{d.3-concept-design-and-publication-status}

Eligible studies meet the immune-link definition in C and report
original human measurements. There is no requirement for the words
``precision,'' ``biomarker,'' ``open window'' or ``recovery,'' for a
particular direction or significance level, or for a
functional/infection endpoint. Studies are neither selected nor weighted
according to whether their authors endorse, reject or ignore the
open-window hypothesis; the author interpretation is charted as a
descriptive tag (C). Designs may include randomized or nonrandomized
experiments, crossover studies, before--after studies and observational
designs that satisfy the timing and comparator rules. No arbitrary
minimum sample size is imposed; small designs are clearly identified and
appraised appropriately.

The primary A/B map includes peer-reviewed original reports. Reviews
support background and citation chasing. Animal-only or in-vitro-only
studies, protocols, editorials, conference abstracts without a full
original report and pure methods papers without eligible human exercise
observations are excluded from the primary map; potentially relevant
abstracts are tracked for full publications.

A bounded supplementary search of bioRxiv and medRxiv will identify
publicly available preprints meeting the same substantive PCC criteria.
They enter a separately labelled frontier register, with separate flow
counts and denominators, and are not pooled with peer-reviewed evidence.
Publication status is checked at the final update, and
preprint/published versions are linked and merged appropriately. Clearly
labelled public preprints may be cited under current Frontiers guidance
{[}14{]}; this does not imply access to unpublished private results.

\subsubsection{D.4 Date, language and
retrieval}\label{d.4-date-language-and-retrieval}

Searches run from database inception to their actual execution dates. No
search language restriction is applied. English and Chinese full texts
are read directly; other languages use translation with independent
human checking of eligibility and key data. Unresolved language or
eligibility information remains awaiting classification rather than
being automatically scientifically excluded.

Institutional access and lawful public copies will be used for full
texts. A sought but unobtained report is ``not retrieved,'' not a
scientific full-text exclusion. This protocol does not plan author
outreach; no contact is sent without separate authorization. Retrieval
attempts and unresolved reports are logged.

\subsection{E. Information sources and search
strategy}\label{e.-information-sources-and-search-strategy}

The five planned databases are PubMed, Web of Science Core Collection,
Scopus, Embase and SPORTDiscus. PubMed is reported as PubMed rather than
implying a separate MEDLINE-platform search. Embase platform choice is
unresolved: draft Embase.com and Ovid translations are supplied for
confirmation; execution of one platform is not claimed as two databases.
The SPORTDiscus draft uses EBSCOhost. Exact collections, platforms,
interfaces and coverage will be recorded.

Two complementary routes are planned: \textbf{E AND I}, where E
represents exercise and I immune concepts; and \textbf{E AND O}, where O
represents omics. Their union is deduplicated. Controlled vocabulary and
free-text fields are adapted to each platform. The immune route covers
immune cells, mediators, function and relevant mucosal/infection terms.
The omics route covers transcriptomics, RNA/non-coding RNA, proteomics,
metabolomics, lipidomics, epigenetics and multi-omics. The exercise
block includes relevant training, competition and sport terminology. The
search is by construct, not by label: ``precision,'' ``biomarker,''
``open window'' and ``recovery'' are not mandatory search concepts, and
the organizing window and all eligibility rules are applied at screening
rather than in the query. The phrase ``open window'' is unsuitable as a
filter---a PubMed title/abstract query combining it with exercise and
immune terms returned 18 records for 2018--2026 at the 4 October 2026
re-check, and original studies rarely use it---so open-window-type terms
may be added only as OR supplements or used for citation chasing, never
as AND terms. Post-exercise and recovery terms remain a diagnostic block
rather than a required concept.

Full working strategies and a search-log template are in
\texttt{03\_search}. They are \textbf{drafts, not platform-validated
searches}. A search specialist or experienced independent search
reviewer will use PRESS {[}11{]}; resulting revisions and resolution of
comments will be retained. Known eligible and boundary examples will
test strategy behavior; the v3 seed list includes open-window-centred
original studies of conventional indices (for example salivary IgA,
natural-killer-cell and lymphocyte redistribution responses after
prolonged exercise) and boundary cases (immediate-only sampling,
sampling only beyond 72 hours, youth cohorts and resting training
studies). Retrieval of a few selected positive seeds is a debugging
check, not an estimate of overall recall. A seed not retrieved will be
investigated; strategy changes must retain a conceptual rationale rather
than simply fitting known papers.

Backward references and forward citations of included reports and
relevant neighboring reviews will be checked, including forward citation
tracking of Peake et al.~{[}15{]} and Campbell and Turner {[}1{]} before
submission. Registered and ongoing reviews in OSF Registries, PROSPERO,
INPLASY, Research Registry and Zenodo, together with adjacent recent
reviews, will inform the novelty assessment; the 4 October 2026 registry
scan is pre-archive reconnaissance and will be repeated before
submission. Supplemental preprint searches will record sites, exact
queries, dates, filters, results examined and stopping rules. Citation
chasing has a defined sweep and documented closure/update; it is not
silently unlimited browsing.

For every run, preserve the exact query, database/platform, run date and
time, limits, hits, export counts, file format, immutable raw export and
checksum. Record search amendments, supplementary routes and
deduplication decisions under PRISMA-S. The target is to update searches
within 30 days before submission; this is a project target, not a
journal rule. New exports do not overwrite earlier exports.

Records are deduplicated using identifiers followed by normalized
title/author/year checks and human confirmation of uncertain matches.
Distinct reports of the same study or cohort remain linked, rather than
being discarded as bibliographic duplicates. Study, report, cohort and
preprint-version relationships are preserved. No formal yields are
available at protocol preparation.

\subsection{F. Selection and
calibration}\label{f.-selection-and-calibration}

Two human reviewers independently screen \textbf{all} titles/abstracts
and \textbf{all} sought full texts. At title/abstract stage, missing
timing or methodological detail normally proceeds to full text rather
than being guessed. AI does not make automatic exclusions or replace
either reviewer.

After the v3.0 protocol has been archived (J) and before formal
screening, a saved random procedure selects 50 records from a frozen
deduplicated draft-search pool; the pool is generated after the archive
release and retained with its version. The pool version, eligible
sampling universe, random seed, selected IDs and reviewer decisions are
retained. The project criterion is at least 80\% raw agreement plus
resolution of all conceptual disagreements; this threshold is not
attributed to JBI. Cohen's kappa is reported descriptively with
prevalence limitations. If the criterion is unmet, the manual is
clarified and a fresh pilot is performed. If the pilot changes an
eligibility rule, the protocol is re-released with justification before
formal searching (J). The 50-record pilot does not confer final study
eligibility. Full-text boundary exercises follow before formal full-text
screening.

Disagreements are discussed; unresolved decisions go to a predesignated
adjudicator whose name must be supplied before formal screening. Changes
to an eligibility interpretation are dated and applied to previously
screened records when relevant. Operational records retain separate
independent decisions and the consensus decision.

The screening manual defines this fixed primary exclusion hierarchy:
FT01 ineligible source type; FT02 not original human research; FT03
ineligible population, including cohorts under 18 years without a
separable adult stratum; FT04 no qualifying identifiable exercise
context, where any intensity, duration or mode qualifies and the code
applies only when no identifiable bout, or for B no qualifying
repeated-bout relation, can be established; FT05 no post-cessation
sample, where sampling only beyond 72 hours is not a reason for this
code because the organizing window is a charting frame; FT06 no baseline
or suitable comparator; FT07 no eligible immune link; FT08 inseparable
mixed intervention. Only the first established scientific reason is
assigned at full-text exclusion; secondary concerns and supporting
passages remain in notes. Unresolved information is not treated as an
established failure. Duplicate reports/versions, not-retrieved reports
and awaiting-classification records have administrative statuses outside
the scientific exclusion hierarchy. The actual flow will reconcile
identification, duplicates, screening, sought/retrieved reports,
exclusions and included reports/cohorts, with the frontier register
separate. Counts remain unavailable, not zero, until work is performed.

\subsection{G. Data charting}\label{g.-data-charting}

\subsubsection{G.1 Pilot, structure and
verification}\label{g.1-pilot-structure-and-verification}

After the archive release, ten purposively chosen reports will calibrate
charting across exercise designs, conventional assays, cell/biofluid
omics, early samples, repeated-bout monitoring, and timing and
population boundary cases, including an adolescent cohort used as a
population boundary case. These reports are training material, not ten
included studies, and may include cases ultimately excluded. Their
source, selection rationale and unresolved questions are recorded in the
pilot manifest. Pilot refinements are versioned before full extraction.

Two humans independently extract eligibility, timing, sample sizes,
effects/uncertainty, precision domains, functional/clinical outcomes,
validation and appraisal. One may enter purely bibliographic metadata
for the other to verify. Discrepancies are resolved against source
passages, with adjudication where necessary. The v3 dictionary, template
and manual in \texttt{05\_extraction} define the operational schema; the
Excel workbooks used for screening, extraction and appraisal are
generated from those JSON templates by
\texttt{scripts/build\_workbooks.py}, and CSV exports are versioned in
git. Data use linked report, study and cohort tables with long-form
assay/marker/time/comparison records. Raw independent extractions and
reconciled data are retained separately.

\subsubsection{G.2 Charting domains}\label{g.2-charting-domains}

\begin{itemize}
\tightlist
\item
  \textbf{Provenance:} Identifiers, publication/version status,
  correction/retraction checks and dates, study/cohort links, source
  locations, extractor and decision history.
\item
  \textbf{Participants:} Eligibility and health definitions, infection
  screening, recruitment/analysis numbers, age range and separability of
  the adult stratum, sex, training status, older-adult subgroup,
  per-group/time/platform participant numbers and within-person pairing.
\item
  \textbf{Exposure:} Design, exercise mode, duration, absolute and
  relative intensity metrics and values as reported, comparator,
  training background, previous load, co-interventions, environmental
  conditions and recovery schedule.
\item
  \textbf{Sampling:} Exact time and time origin, and position relative
  to the organizing window (within window, outside-window recovery,
  pre-exercise baseline or matched control time); matrix/immune
  compartment; collection, processing, storage, delay/freeze--thaw;
  circadian timing; plasma-volume/hydration correction; urinary dilution
  adjustment; salivary flow/secretion rate; cell composition and
  normalization.
\item
  \textbf{Measurements:} Marker family (A); original and standardized
  marker names/identifiers, modality, targeted/untargeted assay,
  platform, units, detection/quantification limits, quality control,
  batch effects, molecular-identification confidence, bulk/single-cell
  structure and integration method.
\item
  \textbf{Results:} Relevant contrasts, direction and magnitude, effect
  estimates and reported CI/SE or other uncertainty, multiplicity
  control, missing data, quality-control exclusions and
  measured/adjusted confounders. Relevant null and discordant findings
  are retained.
\item
  \textbf{Author open-window interpretation:} Whether the report uses
  the open-window term, invokes immunosuppression or susceptibility
  without the term, interprets changes as redistribution, or offers no
  window interpretation, with the source passage; recorded descriptively
  and never used for eligibility or weighting.
\item
  \textbf{Precision domains:} All seven domains in C are charted
  separately: analytical reliability; temporal validity; immune
  specificity; individualization; functional/clinical linkage;
  independent validation and use; and feasibility. Subfields include
  personal baseline, within-person variation, test--retest/repeated-bout
  reliability, heterogeneity analyses, subgroup/individual model
  specification, validation type, participant-level separation,
  calibration/discrimination, thresholds and decision/implementation
  evidence where available.
\item
  \textbf{Functions and outcomes:} Direct assays versus computed
  annotation; co-measurement versus tested association; outcome
  definition/ascertainment, symptom versus clinical versus pathogen
  outcome, biomarker timing relative to onset and follow-up duration.
\end{itemize}

For large omics datasets, chart the assay-level tested universe and all
named immune candidates meeting the original aims/methods definition in
accessible main text and supplements, regardless of significance. Retain
reported feature counts and linked accessible supplements rather than
inventing thousands of unreported rows. Record any candidate-selection
rule and unavailable supplementary data. Preserve result units and avoid
selecting only significant features. Unreported effect sizes are NR.
Figure-derived extraction must identify the source figure/panel, reading
or digitization method, software/settings, units and uncertainty
handling; two reviewers independently check the source and values and
retain their extractions and reconciliation. Estimated values are not
presented as directly reported data. Where reliable extraction is
impossible, use NR or UNCLEAR with the reason.

Participant count, biological sample count, cell count and technical
repeats are different fields. The donor is the independent unit for
single-cell inference; cells do not replace participants. Model splits
are assessed at the participant/cohort level, including feature
selection inside training folds. Raw individual change scores without
measurement/within-person noise assessment do not establish reproducible
responder status.

\subsubsection{G.3 Time representation}\label{g.3-time-representation}

Exact exercise-to-sampling times are primary. Optional display bins are:
pre-exercise baseline; matched control time; during exercise
(ancillary); 0 to \textless30 minutes; 30 minutes to \textless3 hours; 3
to \textless24 hours; 24 to 72 hours inclusive; \textgreater72 hours;
and NR. Bins up to and including 24 to 72 hours lie within the
organizing window; \textgreater72 hours is outside-window recovery. Bins
are mutually exclusive display categories, not eligibility thresholds.
Unknown timing is not imputed to a bin; uncertainty about exercise
linkage may prevent main-map eligibility. Symptom/infection follow-up
uses a separate time axis, with marker measurement before, concurrent
with or after onset recorded explicitly.

\subsection{H. Critical appraisal}\label{h.-critical-appraisal}

Critical appraisal is prespecified to explain evidence limitations, not
to create a clinical recommendation. Two reviewers independently apply
tools matched to actual design, with discussion and adjudication.
Current JBI tools are planned for randomized trials (2023 revision),
quasi-experimental studies (2024), cohorts (2025), analytical
cross-sectional studies (2026), and other eligible designs using the
applicable current checklist {[}12{]}. Single-group pre/post studies are
not treated as randomized trials. Crossover/repeated-measure limitations
are charted explicitly. True diagnostic-accuracy questions use the
current JBI diagnostic-accuracy tool. True individual-level future
clinical-outcome or immune-response prediction-model
development/validation studies additionally use PROBAST+AI (2025)
{[}13{]}; a statistical model, molecular discovery classifier or simple
association alone does not trigger that tool.

The exact official forms, versions, access dates and checksums will be
archived before appraisal. Native response categories, domains and
supporting source passages are retained. We do not assume all JBI forms
have identical item structure or force them into a common score. The
manual in \texttt{05\_extraction/critical\_appraisal\_manual.md} governs
design selection and overlap.

Assay/omics reporting and the seven precision domains are recorded
separately from design appraisal. They include analytical reliability,
pre-analytical handling, normalization, cell composition,
repeated-measure analysis, multiplicity and independent validation.
These custom descriptors are not a validated risk-of-bias tool. There is
no quality sum, quality-threshold exclusion, GRADE certainty assessment
or numerical readiness ranking; domain evidence enters only the
validation-readiness map as defined in A. Appraisal findings
contextualize the maps, including unresolved reporting limitations.

\subsection{I. Synthesis and planned
outputs}\label{i.-synthesis-and-planned-outputs}

Synthesis is descriptive, stratified and narrative. Independent cohorts
are the primary count for evidence density; report counts and distinct
studies are also reported with explicit linkage. Platform features,
cells, assays or time points are never used as proxies for evidence
strength. Denominators are specific to each question and include
applicable/NR/UNCLEAR information as appropriate. A/B overlap is
disclosed, and the two stream totals are not simply added.

Three primary evidence maps are planned:

\begin{enumerate}
\def\labelenumi{\arabic{enumi}.}
\tightlist
\item
  \textbf{Open-window time by immune compartment and marker family:}
  Exact sampling coverage and optional time bins within and beyond the
  organizing window, by immune compartment or matrix, marker family,
  exercise context and population stratum.
\item
  \textbf{Per-marker validation-readiness map:} Marker or marker family
  by the seven precision domains, with conventional and omics candidates
  distinguishable; the descriptive profile defined in A, not a score or
  ranking.
\item
  \textbf{Cohort-level validation links:} Technical confirmation,
  repeated bouts, independent participants/cohorts, directly measured
  function, tested associations and clinical/prediction outcomes, with
  explicitly labelled gaps.
\end{enumerate}

The PRISMA flow is an additional process figure. Main tables describe
study/cohort context, measurement/timing and limitations,
individualization/validation evidence, and differences from adjacent
reviews. Core A, support B and the frontier preprint register are
visibly separated. Background sources do not enter original-study
denominators. Counts of studies reporting a direction may be descriptive
but are not a vote-counting estimate of effect or truth. Null,
contradictory and unreported findings remain visible.

We will not pool effects, derive a diagnostic threshold, validate a
marker, recommend return-to-training decisions, conclude that a
candidate predicts infection from group differences alone, or adjudicate
whether an open window exists. Discovery of a potentially homogeneous
subset will be documented as a future focused-review opportunity. Any
later decision to add an effect synthesis requires an explicit protocol
amendment and appropriate methods before analysis. The manuscript will
retain the scoping methodology regardless of the submission portal's
article-category label {[}14{]}.

\subsection{J. Registration, archiving, governance and
transparency}\label{j.-registration-archiving-governance-and-transparency}

The protocol is archived publicly rather than registered in a review
registry. After the research team freezes version 3.0, the public GitHub
repository will be tagged and a GitHub Release archived automatically by
Zenodo, which assigns a dated, immutable version DOI and a concept DOI
for the record series; the DOI is assigned on release and cannot be
stated earlier. The release precedes the 50-record screening pilot and
the 10-report charting pilot, and formal five-database searching starts
only after the archived release exists. If the pilots change a rule,
version 3.1 is released with justification, as a new tag and version DOI
under the same concept DOI, before formal searching. PROSPERO does not
accept scoping reviews, and OSF Registries is no longer the planned
route. At the date of this draft no release, DOI or registration of any
kind exists, and the project has never been registered in any registry.

Pre-archive knowledge is disclosed rather than denied. An internal
protocol (version 1.0, January 2026) with its own literature collision
assessments, public reconnaissance and source verification on 2 October
2026 (Codex), and a collision re-check of the narrowed framing against
published and registered reviews on 4 October 2026 (Claude Code;
\texttt{02\_preliminary/collision\_recheck\_2026-10-04.md}) all precede
the archive, and the two pilots are planned stages that will also
precede formal searching. Their dates, scope and influence on decisions
are recorded in the archive release record and the amendment log. The
work is not described as wholly prospective from before any knowledge of
the literature.

Substantive amendments record the previous rule, new rule, reason, date,
work stage, information already known, authorizing team decision and
need to revisit earlier records; after the archive release, each
amendment also produces a new tagged Release and version DOI. Raw search
exports and independent decisions remain immutable. The canonical
protocol and execution appendices are versioned together.

This review uses published or lawfully accessible material and does not
recruit participants or intervene in exercise. Institutional
requirements for ethical review/exemption will be checked and recorded;
no approval or exemption number is invented. The GitHub repository and
Zenodo archive hold the protocol, manuals, templates, search strategies,
dictionaries, decision metadata and shareable synthesis data under CC BY
4.0 for documents and MIT for scripts; copyrighted full texts and raw
database exports are not redistributed contrary to their terms. Author
names, affiliations, funding and competing interests are pending human
declarations.

AI has supported public-source reconnaissance, evidence organization,
consistency checking and drafting in this project: Codex desktop with
delegated agents on 2--3 October 2026, and Claude Code (Anthropic) on 4
October 2026 for the collision re-check and this v3.0 revision. Tool
names, versions, models, sources, dates, tasks, prompt/output references
and human verification are logged in
\texttt{01\_protocol/ai\_use\_log.json}, and the manuscript will
disclose tool name, version, model and source in both Methods and
Acknowledgments {[}14{]}. AI is not an author, human independent
reviewer or adjudicator. Named humans remain responsible for
eligibility, data, appraisal, citations and conclusions. No AI-produced
citation or numerical claim is accepted without source verification.
Final disclosure follows the then-current journal and institutional
requirements.

\subsection{K. Roles and execution
plan}\label{k.-roles-and-execution-plan}

Roles to name are: principal investigator; reviewers A and B; search
lead; independent PRESS reviewer; adjudicator; data/version manager; and
analysis/writing lead. A second human reviewer and access to all five
databases/full texts are available. Other roles may overlap where
independence is preserved, but an adjudicator must be named before
formal screening. Funding and competing interests cannot be inferred
from the absence of information.

Screening, extraction and appraisal use Excel workbooks generated from
the JSON templates in \texttt{04\_screening} and \texttt{05\_extraction}
by \texttt{scripts/build\_workbooks.py}: a read-only record master, one
locked decision workbook per reviewer, dropdown-controlled codes, and
identifier and gene-name columns stored as text. Each reviewer's
workbook is locked at completion and its SHA-256 hash committed to git
before the two files are merged; \texttt{scripts/merge\_screening.py}
performs the deduplication checks, merge, agreement, kappa and PRISMA
counts, and CSV exports are versioned in git. Rayyan, Covidence or
similar platforms are not used. A reference manager with exportable
identifiers is still to be chosen; software and versions are recorded
before execution.

Preparation is estimated at approximately 2--3 working weeks, depending
on team availability: 3--5 working days to verify platforms and search
strategies and to generate the workbooks; 2--3 working days to resolve
team comments, freeze version 3.0 and publish the archived release; and
5--7 working days for reviewer training and the pilot screening and
charting, with parallel work where feasible. If the pilots change a
rule, a conditional re-release follows before formal searching. These
are planning estimates, not a promised completion date. Full-review
effort will be estimated after the 50-record screening pilot and
10-report charting pilot, using actual yields and per-record effort. A
provisional budget should distinguish two-person screening/extraction
from single-person work and allow adjudication and update searches.

The local execution plan and archive release record specify stage gates
and deliverables. Formal stages remain pending until the required human
work is completed; completing this document does not certify PRESS
review, calibration or archiving.

\subsection{L. Version history and decision
status}\label{l.-version-history-and-decision-status}

Version 1.0 (January 2026) and associated build artifacts are preserved
in the dated archive. Version 2.0 (2 October 2026) replaced the
open-window-led title with candidate biomarkers for post-exercise
precision monitoring; defined seven independent precision domains;
separated acute A from strictly time-linked repeated-bout B; retained
immediate samples without a 72-hour ceiling; expanded searching to
inception without a language filter; removed mandatory
precision/biomarker/open-window terms; added cohort linkage, complete
outcome/validation fields, design-appropriate descriptive appraisal and
AI transparency; and removed unsupported novelty claims.

Version 3.0 (this draft, 4 October 2026) follows the collision re-check
of the same date. It re-adopts the open window as the organizing
construct that defines the 0--72 hour post-exercise window, the sampling
time points and the marker families, without adjudicating whether it
exists; defines the per-marker validation-readiness map as the
deliverable; removes any exercise intensity, duration or mode threshold;
restricts the population to adults aged 18 years or older, with older
adults as a subgroup; retains Support B as the cross-bout monitoring
layer; adds the author open-window interpretation tag, marker family and
window-position fields; replaces OSF registration with a public GitHub
repository and a Zenodo-archived Release, to be published after team
freeze and before the pilots; specifies Excel workbooks generated from
the JSON templates with scripted merging and agreement; adopts CC BY 4.0
for documents and MIT for scripts; and adds the Frontiers in Immunology
submission requirements (PRISMA-ScR checklist and flow diagram, AI-use
disclosure in Methods and Acknowledgments). The v3.0 date becomes the
actual team freeze date at release.

The scope and operational design are ready for team freeze. Still
pending are named administrative declarations, platform-specific
verification/PRESS review, the Zenodo release, human pilots, archived
appraisal forms and local software configuration. Subsequent formal
searches, screening, extraction and synthesis are distinct execution
stages and have not been performed.

\subsection{References}\label{references}

\begin{enumerate}
\def\labelenumi{\arabic{enumi}.}
\tightlist
\item
  Campbell JP, Turner JE. Debunking the Myth of Exercise-Induced Immune
  Suppression: Redefining the Impact of Exercise on Immunological Health
  Across the Lifespan. \emph{Frontiers in Immunology}. 2018.
  \href{https://doi.org/10.3389/fimmu.2018.00648}{doi:10.3389/fimmu.2018.00648}.
\item
  Davison G, et al.~Rigorous methodological approaches to address
  knowledge gaps in exercise, nutrition, immunity, and infection risk
  research. \emph{Physiological Reports}. 2025.
  \href{https://doi.org/10.14814/phy2.70479}{doi:10.14814/phy2.70479}.
\item
  FDA--NIH Biomarker Working Group. BEST (Biomarkers, EndpointS, and
  other Tools) Resource: Validation.
  \href{https://www.ncbi.nlm.nih.gov/books/NBK464453/}{NCBI Bookshelf}.
  Accessed 2 October 2026.
\item
  Kotewitsch M, et al.~Non-coding RNAs in exercise immunology: A
  systematic review. \emph{Journal of Sport and Health Science}. 2024.
  \href{https://doi.org/10.1016/j.jshs.2023.11.001}{doi:10.1016/j.jshs.2023.11.001}.
\item
  Walzik D, et al.~Impact of exercise on markers of B cell-related
  immunity: A systematic review. \emph{Journal of Sport and Health
  Science}. 2024.
  \href{https://doi.org/10.1016/j.jshs.2023.10.002}{doi:10.1016/j.jshs.2023.10.002}.
\item
  Teodoro TH, et al.~The effects of acute and chronic exercise on immune
  markers of TH1/TH2 cells in older adults: a systematic review.
  \emph{Frontiers in Physiology}. 2025.
  \href{https://doi.org/10.3389/fphys.2025.1453747}{doi:10.3389/fphys.2025.1453747}.
\item
  Minuzzi LG, et al.~Immunometabolomics Applied to Physical Exercise:
  Accomplishments and New Directions for Health Improvement.
  \emph{Annual Review of Analytical Chemistry}. 2026.
  \href{https://doi.org/10.1146/annurev-anchem-091024-095826}{doi:10.1146/annurev-anchem-091024-095826}.
\item
  Pollock D, Peters MDJ, Tricco AC, et al.~Scoping reviews (2026). In:
  Aromataris E, et al., editors. \emph{JBI Manual for Evidence
  Synthesis}. JBI; 2024.
  \href{https://jbi-global.atlassian.net/wiki/spaces/MANUAL/pages/355862497}{Current
  online chapter}.
  \href{https://doi.org/10.46658/JBIMES-24-09}{doi:10.46658/JBIMES-24-09}.
  Accessed 2 October 2026.
\item
  Tricco AC, et al.~PRISMA Extension for Scoping Reviews (PRISMA-ScR):
  Checklist and Explanation. \emph{Annals of Internal Medicine}. 2018.
  \href{https://doi.org/10.7326/M18-0850}{doi:10.7326/M18-0850}.
\item
  Rethlefsen ML, et al.~PRISMA-S: an extension to the PRISMA Statement
  for Reporting Literature Searches in Systematic Reviews.
  \emph{Systematic Reviews}. 2021.
  \href{https://doi.org/10.1186/s13643-020-01542-z}{doi:10.1186/s13643-020-01542-z}.
\item
  McGowan J, et al.~PRESS Peer Review of Electronic Search Strategies:
  2015 Guideline Statement. \emph{Journal of Clinical Epidemiology}.
  2016.
  \href{https://doi.org/10.1016/j.jclinepi.2016.01.021}{doi:10.1016/j.jclinepi.2016.01.021}.
\item
  JBI. \href{https://jbi.global/critical-appraisal-tools}{Critical
  appraisal tools}. Current design-specific forms; accessed 2 October
  2026.
\item
  Moons KGM, et al.~PROBAST+AI: an updated quality, risk of bias, and
  applicability assessment tool for prediction models using regression
  or artificial intelligence methods. \emph{BMJ}. 2025;388:e082505.
  \href{https://doi.org/10.1136/bmj-2024-082505}{doi:10.1136/bmj-2024-082505}.
\item
  Frontiers.
  \href{https://www.frontiersin.org/journals/immunology/for-authors/article-types}{Immunology
  article types}, including the Systematic Review type and its PRISMA
  reporting and flow-diagram requirement;
  \href{https://www.frontiersin.org/guidelines/author-guidelines}{author
  guidelines}; and
  \href{https://www.frontiersin.org/guidelines/policies-and-publication-ethics\#artificial-intelligence-policy}{policies
  and publication ethics}, including reporting guidelines and the
  artificial-intelligence fair use and disclosure policy. Accessed 4
  October 2026; recheck before submission.
\item
  Peake JM, Neubauer O, Walsh NP, Simpson RJ. Recovery of the immune
  system after exercise. \emph{Journal of Applied Physiology}.
  2017;122(5):1077--1087.
  \href{https://doi.org/10.1152/japplphysiol.00622.2016}{doi:10.1152/japplphysiol.00622.2016}.
  PMID 27909225.
\item
  Simpson RJ, Campbell JP, Gleeson M, et al.~Can exercise affect immune
  function to increase susceptibility to infection? \emph{Exercise
  Immunology Review}. 2020;26:8--22. PMID 32139352; no DOI.
\item
  Xie E, He Y, Cao T, et al.~Exercise and tissue-resident memory T
  cells: from circulating numbers to spatial immune remodeling.
  \emph{Frontiers in Immunology}. 2026;17:1822561.
  \href{https://doi.org/10.3389/fimmu.2026.1822561}{doi:10.3389/fimmu.2026.1822561}.
  PMID 42183211.
\item
  Shi X, Hu L, Nieman DC, et al.~Exercise workload: a key determinant of
  immune health -- a narrative review. \emph{Frontiers in Immunology}.
  2025;16:1617261.
  \href{https://doi.org/10.3389/fimmu.2025.1617261}{doi:10.3389/fimmu.2025.1617261}.
  PMID 40777031.
\item
  Reitzner SM, Brodin P. The Immune System as a Biosensor of Health and
  Athletic Performance. \emph{Scandinavian Journal of Medicine \&
  Science in Sports}. 2026;36(3):e70258.
  \href{https://doi.org/10.1111/sms.70258}{doi:10.1111/sms.70258}. PMID
  41842715.
\item
  Mănescu DC, Voinea A, Plastoi CD, et al.~Molecular Biomarkers of
  Training Responses: A Systems Framework for Exercise Adaptation and
  Athlete Monitoring. \emph{International Journal of Molecular
  Sciences}. 2026;27(8):3601.
  \href{https://doi.org/10.3390/ijms27083601}{doi:10.3390/ijms27083601}.
  PMID 42074239.
\end{enumerate}

\end{document}
